\documentclass[10pt]{beamer}
\usetheme{Madrid}
\usepackage{booktabs}
\usepackage{graphicx}
\usepackage[T1]{fontenc}
\setbeamertemplate{navigation symbols}{}

% ======================= EDIT YOUR DETAILS =======================
\newcommand{\studentname}{Aflah Muhammed P}
% Change the university roll/registration number:
\newcommand{\rollno}{TVE23CSXXX}
% Change your class roll number:
\newcommand{\rollclass}{Roll No: XX}
% Change your guide name:
\newcommand{\guidename}{Prof. Guide Name}
% Change department / college / short-institute if needed:
\newcommand{\deptname}{Department of Computer Science and Engineering}
\newcommand{\collegename}{College of Engineering Trivandrum}
\newcommand{\shortinst}{CET}
% Change the seminar date:
\newcommand{\seminardate}{\today}
% =================================================================

\title[B.Tech Seminar]{How AI Agents Team Up: Understanding Multi-Agent Collaboration in Large Language Models}
\subtitle{Based on: Multi-Agent Collaboration Mechanisms: A Survey of LLMs (arXiv:2501.06322, 2025)}
\author[\studentname\ (\shortinst)]{\studentname \\ \small \rollno \\ \small \rollclass \\ \small Guide: \guidename}
\institute[]{\deptname \\ \collegename}
\date{\seminardate}

\begin{document}

\begin{frame}[plain]
  \titlepage
\end{frame}

\begin{frame}{Table of Contents}
  \tableofcontents
\end{frame}

\section{Introduction}
\begin{frame}{Introduction}
\begin{itemize}
  \item AI is moving from single models to teams of cooperating LLM agents
  \item Each agent has a goal, memory, and can act and communicate
  \item Agent teams tackle complex, multi-step tasks better than one model
  \item The field grew fast but lacked shared vocabulary and structure
  \item This survey maps how LLM agents collaborate into one framework
\end{itemize}
\end{frame}

\section{Literature Survey}
\begin{frame}{Literature Survey / Background}
  \small
  \begin{tabular}{@{}p{2.7cm}p{2.3cm}p{5.6cm}@{}}
    \toprule
    \textbf{Title} & \textbf{Author} & \textbf{Summary} \\
    \midrule
    CAMEL: Communicative Agents for Mind Exploration & Guohao Li et al., 2023 & Role-playing framework where two cooperating agents converse through assigned roles to complete a shared task autonomously. \\
    \addlinespace
    MetaGPT: Meta Programming for Multi-Agent Collaboration & Sirui Hong et al., 2024 & Encodes human Standard Operating Procedures so role-based agents work assembly-line style to build software collaboratively. \\
    \addlinespace
    Improving Factuality and Reasoning through Multiagent Debate & Yilun Du et al., 2023 & Multiple agents debate and critique each other's answers, improving factual accuracy and reasoning over a single model. \\
    \addlinespace
    \bottomrule
  \end{tabular}
\end{frame}

\section{Motivation}
\begin{frame}{Motivation}
\begin{itemize}
  \item Single LLMs struggle with large, complex, multi-step tasks
  \item Many real problems split naturally into specialized roles
  \item Rapidly growing MAS research had no unifying framework
  \item Hard to compare, reuse, or secure systems without shared terms
\end{itemize}
\end{frame}

\section{Problem Statement}
\begin{frame}{Problem Statement}
  \begin{block}{}
  There is no common framework or vocabulary for describing how LLM-based agents collaborate, which makes multi-agent systems hard to compare, design, and improve. This survey aims to characterize and organize those collaboration mechanisms.
  \end{block}
\end{frame}

\section{Objectives}
\begin{frame}{Objectives}
\begin{itemize}
  \item Survey the collaborative aspect of LLM multi-agent systems
  \item Introduce an extensible framework with clear key dimensions
  \item Classify existing systems using consistent shared vocabulary
  \item Review applications across networks, industry, QA, and society
  \item Identify open challenges and future research directions
\end{itemize}
\end{frame}

\section{Methodology}
\begin{frame}[allowframebreaks]{Methodology / Approach}
\begin{itemize}
  \item Define a single agent by model, objective, environment, input, output
  \item Characterize collaboration along five key dimensions
  \item Actors: which agents join and what expertise they bring
  \item Types: cooperation, competition, or coopetition
  \item Structures: centralized, decentralized, or hierarchical
  \item Strategies and protocols: rule/role/model-based, static or dynamic
\end{itemize}
\end{frame}

\section{Key Concepts}
\begin{frame}[allowframebreaks]{Key Concepts}
  \begin{description}
    \item[Actors] The agents involved, each assigned objectives based on their expertise and role in the team.
    \item[Interaction types] Cooperation toward shared goals, competition for individual goals, or coopetition mixing both.
    \item[Structures] Team shape: centralized hub, decentralized peer-to-peer, or layered hierarchical organization.
    \item[Strategies] How behavior is controlled: fixed rules, specialized roles, or reasoning about others' intentions.
    \item[Coordination protocols] Static preset communication channels, or dynamic ones adjusted by a manager or orchestration graph.
  \end{description}
\end{frame}

\section{System Architecture}
\begin{frame}{System Architecture / How It Works}
  \begin{block}{Overview}
  Draw each agent as a box holding its model (LLM plus memory and system prompt), its objective, and its input-to-output action loop. Then show a group of these agents connected by message arrows, with the connection pattern set by the chosen structure: a centralized layout has one hub agent routing all messages, a decentralized layout has agents linked peer-to-peer, and a hierarchical layout stacks agents in layers of authority. A coordination protocol governs the turn order and handoffs along those arrows, either fixed in advance (static) or adjusted on the fly by a manager agent or an orchestration graph (dynamic). Labeling the interaction type (cooperation, competition, or coopetition) on the whole diagram completes the picture.
  \end{block}
  % TIP: insert a diagram here, e.g.:
  % \begin{center}\includegraphics[width=0.7\textwidth]{your-figure.png}\end{center}
\end{frame}

\section{Results}
\begin{frame}[allowframebreaks]{Results / Findings}
\begin{itemize}
  \item Delivers a unified five-dimension collaboration framework
  \item Classifies systems like CAMEL, AutoGen, MetaGPT, ChatDev consistently
  \item Organizes three interaction types and three structures
  \item Covers four domains: 5G/6G, Industry 5.0, QA, social/cultural
  \item Notes multi-agent debate tends to improve factuality and reasoning
\end{itemize}
\end{frame}

\section{Discussion}
\begin{frame}{Discussion}
\begin{itemize}
  \item Shared vocabulary lets teams compare and reuse designs
  \item Choice of structure and strategy should match the task
  \item Agent teams broaden LLM impact across many real domains
  \item Framework is a practical design checklist, not just theory
\end{itemize}
\end{frame}

\section{Challenges}
\begin{frame}{Limitations \& Challenges}
\begin{itemize}
  \item LLMs are standalone models, not trained for collaboration
  \item Agent behavior is hard to explain and predict
  \item Coordination cost and complexity grow with more agents
  \item Safety risks: hallucination spread and cascading failures
\end{itemize}
\end{frame}

\section{Future Work}
\begin{frame}{Future Work}
\begin{itemize}
  \item Design better collaboration channels and adaptive role assignment
  \item Integrate domain knowledge and pick strategies per task
  \item Add safety and ethical governance for agent teams
  \item Advance toward scalable artificial collective intelligence
\end{itemize}
\end{frame}

\section{Conclusion}
\begin{frame}{Conclusion}
\begin{itemize}
  \item LLM agents are shifting AI from lone models to teams
  \item A five-dimension framework maps how these teams collaborate
  \item It classifies existing systems and real-world applications
  \item Open challenges in scale, cost, and safety guide future work
\end{itemize}
\end{frame}

\section{References}
\begin{frame}[allowframebreaks]{References}
  \footnotesize
  \begin{enumerate}
    \item Multi-Agent Collaboration Mechanisms: A Survey of LLMs, Khanh-Tung Tran, Dung Dao, Minh-Duong Nguyen, Quoc-Viet Pham, Barry O'Sullivan, Hoang D. Nguyen, arXiv:2501.06322, 2025
    \item CAMEL: Communicative Agents for Mind Exploration of Large Language Model Society, Guohao Li et al., NeurIPS, 2023
    \item AutoGen: Enabling Next-Gen LLM Applications via Multi-Agent Conversation, Qingyun Wu et al., arXiv:2308.08155, 2023
    \item MetaGPT: Meta Programming for a Multi-Agent Collaborative Framework, Sirui Hong et al., ICLR, 2024
    \item Communicative Agents for Software Development (ChatDev), Chen Qian et al., ACL, 2024
    \item Improving Factuality and Reasoning in Language Models through Multiagent Debate, Yilun Du et al., ICML, 2023
  \end{enumerate}
\end{frame}

\begin{frame}[plain]
  \begin{center}
    {\Huge Thank You}\\[1em]
    {\large Questions?}
  \end{center}
\end{frame}

\end{document}